\begin{proof}[Proof of \cref{obj:lem:dimension-free-private-two-point}]
Set
\[
c_1=\frac{1}{8192}>0.
\]
Fix admissible \(n,d,\varepsilon\) and a sampling scheme satisfying
\cref{obj:ass:iid-people,obj:ass:independent-randomness}.

\begin{enumerate}
\item Fix \(\mathcal C\in\{\mathrm{NI},\mathrm{SI}\}\) and an admissible protocol
\[
Q\in\mathfrak Q_{\mathcal C}(n,d,\varepsilon).
\]
If \(\mathcal C=\mathrm{NI}\), the inclusion in
\cref{obj:def:noninteractive-class} gives
\(Q\in\mathfrak Q_{\mathrm{SI}}(n,d,\varepsilon)\).
If \(\mathcal C=\mathrm{SI}\), this membership is the stipulated admissibility condition. Thus in both cases,
\[
Q\in\mathfrak Q_{\mathrm{SI}}(n,d,\varepsilon).
\]
We establish bounds uniform over all decisions based on this protocol.
% lean: CausalSmith.Stat.LdpOptvalueUniformFrontier.dimension_free_private_two_point

\item Define the positive scale and two-point amplitude by
\[
\kappa_{\mathrm{tp}}=\sqrt n\,\varepsilon>0,
\qquad
\alpha_{\mathrm{tp}}
=\frac18\min\left\{1,\frac1{\kappa_{\mathrm{tp}}}\right\}.
\]
Since \(n\ge2\) and \(\varepsilon>0\),
\[
0<\alpha_{\mathrm{tp}}\le\frac18\le\frac12,
\qquad
\alpha_{\mathrm{tp}}\sqrt n\,\varepsilon
=\frac{\kappa_{\mathrm{tp}}}{8}
  \min\left\{1,\frac1{\kappa_{\mathrm{tp}}}\right\}
\le\frac18.
\]
For the squared-rate calibration, splitting according to whether
\(\kappa_{\mathrm{tp}}\le1\) or \(\kappa_{\mathrm{tp}}>1\) gives
\[
\left(\min\left\{1,\frac1{\kappa_{\mathrm{tp}}}\right\}\right)^2
=
\begin{cases}
1,&\kappa_{\mathrm{tp}}\le1,\\
\kappa_{\mathrm{tp}}^{-2},&\kappa_{\mathrm{tp}}>1
\end{cases}
=
\min\left\{1,\frac1{\kappa_{\mathrm{tp}}^2}\right\}.
\]
Here the first case uses
\(1/\kappa_{\mathrm{tp}}\ge1\) and
\(1/\kappa_{\mathrm{tp}}^2\ge1\), while the second uses the reverse inequalities.
As
\[
\kappa_{\mathrm{tp}}^2=n\varepsilon^2,
\]
we obtain the two numerical comparisons
\[
\frac{63\alpha_{\mathrm{tp}}^2}{4096}
=
\frac{63}{262144}\min\left\{1,\frac1{n\varepsilon^2}\right\}
\ge
c_1\min\left\{1,\frac1{n\varepsilon^2}\right\},
\]
and
\[
\frac{81\alpha_{\mathrm{tp}}}{320}
=
\frac{81}{2560}\min\left\{1,\frac1{\sqrt n\,\varepsilon}\right\}
\ge
c_1\min\left\{1,\frac1{\sqrt n\,\varepsilon}\right\}.
\]
% lean: CausalSmith.Stat.LdpOptvalueUniformFrontier.two_point_amplitude_calibration

\item Define the two contrast vectors, coordinate priors, and privacy contraction scale by
\[
\theta^{(0)}_j=0,\qquad
\theta^{(1)}_j=\alpha_{\mathrm{tp}}
\quad(j\in[d]),
\qquad
\nu_0=\operatorname{Dirac}_0,\qquad
\nu_1=\operatorname{Dirac}_{\alpha_{\mathrm{tp}}},
\qquad
\beta=\frac{e^\varepsilon-1}{2d},
\]
where \(\operatorname{Dirac}_x\) denotes the unit point mass at \(x\).
These priors satisfy
\[
\nu_0([-\alpha_{\mathrm{tp}},\alpha_{\mathrm{tp}}])
=\nu_1([-\alpha_{\mathrm{tp}},\alpha_{\mathrm{tp}}])=1,
\qquad
\int u^0\,\nu_0(du)=\int u^0\,\nu_1(du)=1.
\]
Thus their moments match through order zero. Each coordinate under the corresponding product prior equals its point-mass location almost surely; intersecting these finitely many probability-one events yields
\[
\nu_\ell^{\otimes d}\bigl(\{\theta^{(\ell)}\}\bigr)=1
\qquad(\ell\in\{0,1\}).
\]
Consequently, writing the two joint seed-transcript laws as
\[
N_\ell=\mathcal L_{G_{\theta^{(\ell)}},Q}(R,\mathbf Z)
\qquad(\ell\in\{0,1\}),
\]
the mixtures in \cref{obj:lem:adaptive-moment-contraction} reduce to
\[
\mathsf M_\ell^Q=N_\ell
\qquad(\ell\in\{0,1\}).
\]
Apply that result with amplitude \(\alpha_{\mathrm{tp}}\), these coordinate priors, and \(k=1\). Its hypotheses hold by the assumed measurable Radon--Nikodym property, the sampling conditions, sequential membership of \(Q\), and \(1\le n\). It gives
\[
\begin{aligned}
\operatorname{TV}(N_0,N_1)
&\le
\min\left\{1,\,
d\alpha_{\mathrm{tp}}\sqrt{\binom n1}\,\beta\right\}\\
&\le
\frac{\alpha_{\mathrm{tp}}\sqrt n\,(e^\varepsilon-1)}2.
\end{aligned}
\]
For \(0<\varepsilon\le1\), the standard exponential estimate gives
\[
e^\varepsilon-1
\le |e^\varepsilon-1|
\le2|\varepsilon|
=2\varepsilon.
\]
Therefore,
\[
\operatorname{TV}(N_0,N_1)
\le\alpha_{\mathrm{tp}}\sqrt n\,\varepsilon
\le\frac18.
\]
% lean: CausalSmith.Stat.LdpOptvalueUniformFrontier.constant_contrast_tv_bound

\item Both contrast vectors lie in \([-1/2,1/2]^d\). By
\cref{obj:prop:signed-causal-bridge}, the symmetric laws of
\cref{obj:def:symmetric-law} satisfy
\[
G_{\theta^{(0)}},G_{\theta^{(1)}}\in\mathcal V_d.
\]
Since \(d>0\), their average absolute contrasts are
\[
F(\theta^{(0)})=\frac1d\sum_{j=1}^d0=0,
\qquad
F(\theta^{(1)})
=\frac1d\sum_{j=1}^d\alpha_{\mathrm{tp}}
=\alpha_{\mathrm{tp}}.
\]
The same cited result therefore gives
\[
V(G_{\theta^{(0)}})=\frac12,
\qquad
V(G_{\theta^{(1)}})=\frac12+\frac{\alpha_{\mathrm{tp}}}{2}.
\]
Define the full-data priors and their welfare centers by
\[
\Pi_0=\operatorname{Dirac}_{G_{\theta^{(0)}}},
\qquad
\Pi_1=\operatorname{Dirac}_{G_{\theta^{(1)}}},
\qquad
v_0=\frac12,
\qquad
v_1=\frac12+\frac{\alpha_{\mathrm{tp}}}{2},
\qquad
\Delta=v_1-v_0=\frac{\alpha_{\mathrm{tp}}}{2}>0.
\]
Their causal support and concentration properties are
\[
\Pi_\ell(\mathcal V_d)=1,
\qquad
\Pi_\ell\left(\left\{P:
|V(P)-v_\ell|>\frac{\Delta}{8}\right\}\right)=0
\qquad(\ell\in\{0,1\}).
\]
Mixing over a unit point mass returns the corresponding decision law. These decision laws are probability measures because the participant and randomization laws are probability measures and the message kernels are Markov kernels. Their seed-transcript marginals are precisely \(N_0,N_1\).

Set
\[
\eta_0=0,\qquad \omega=\frac18.
\]
The preceding displays discharge the hypotheses of
\cref{obj:lem:separated-value-mixtures}. Hence every admissible estimator \(T\) satisfies
\[
\sup_{P\in\mathcal V_d}
\mathbb E_{P,Q}\bigl[(T(\mathbf Z,R,U)-V(P))^2\bigr]
\ge
\frac{9\Delta^2}{128}\cdot\frac78
=
\frac{63\alpha_{\mathrm{tp}}^2}{4096}.
\]
Every admissible connected interval \(I\) with global coverage \(0.90\) satisfies
\[
\sup_{P\in\mathcal V_d}
\mathbb E_{P,Q}\operatorname{length}I(\mathbf Z,R,U)
\ge
\frac{3\Delta}{4}\cdot\frac{27}{40}
=
\frac{81\alpha_{\mathrm{tp}}}{320}.
\]
Here \(7/8=1-\omega-2\eta_0\) and
\(27/40=0.80-2\eta_0-\omega\), so both positive-part factors equal the displayed positive numbers.
% lean: CausalSmith.Stat.LdpOptvalueUniformFrontier.constant_contrast_decision_lower

\item Combining these bounds with the amplitude calibration yields
\[
\sup_{P\in\mathcal V_d}
\mathbb E_{P,Q}\bigl[(T(\mathbf Z,R,U)-V(P))^2\bigr]
\ge
c_1\min\left\{1,\frac1{n\varepsilon^2}\right\}
\]
for every admissible estimator, and
\[
\sup_{P\in\mathcal V_d}
\mathbb E_{P,Q}\operatorname{length}I(\mathbf Z,R,U)
\ge
c_1\min\left\{1,\frac1{\sqrt n\,\varepsilon}\right\}
\]
for every admissible globally honest interval.
% lean: CausalSmith.Stat.LdpOptvalueUniformFrontier.dimension_free_private_two_point_decisions

\item These inequalities hold for every admissible protocol in either selected class. Taking the infima over protocols and decisions in
\cref{obj:def:risk,obj:def:honest-length} therefore gives
\[
\mathfrak R_{\mathcal C}(n,d,\varepsilon)
\ge
c_1\min\left\{1,\frac1{n\varepsilon^2}\right\},
\qquad
\mathfrak H_{\mathcal C}(n,d,\varepsilon)
\ge
c_1\min\left\{1,\frac1{\sqrt n\,\varepsilon}\right\}.
\]
The infima are interpreted in the extended nonnegative reals, whose infimum over an empty decision class is \(+\infty\). Thus the same lower bounds cover that case as well.
% lean: CausalSmith.Stat.LdpOptvalueUniformFrontier.dimension_free_private_two_point
\end{enumerate}
\end{proof}
